This paper presented GaitNet, a greedy footstep planner for acyclic
quadruped locomotion. GaitNet makes one decision every
\corlControlTickMs\,ms. It can select one terrain-valid footstep for
one leg, choose the swing duration, or take no action.
\changed{At deployment, the selected footstep is the exact maximizer
of GaitNet's learned value function over every valid candidate cell
and the learned no-action option.} Each swing lasts longer than one
planning tick. As a result, overlapping swings can emerge from
repeated decisions. The planner does not represent a gait phase, leg
pair, or contact sequence.

\pending{Hardware trials on a Unitree Go1 show that GaitNet (Full)
outperforms the Legged Perceptive Stack baseline. It also outperforms
the concurrency-restricted GaitNet-Single and GaitNet-Trot ablations.
GaitNet (Full) achieves a \textbf{\corlHwFullSuccess\%} mean success
rate, compared with \textbf{\corlHwLPSSuccess\%} for the scheduled
NMPC baseline. The advantage is largest at the highest tested terrain
difficulties.} Preliminary simulation results from our earlier
framework \cite{Sullivan2025} show the same qualitative mechanism.
The policy steps one foot at a time under easy commands. Under harder
commands, it produces overlapping, non-periodic swings with no change
in the decision rule. \pending{The hardware results indicate that
this high-rate re-decision mechanism transfers to physical
deployment.}

\subsection{Limitations and Future Work}

GaitNet's concurrency is still limited by an explicit feasibility
constraint. At least \numberstringnum{\corlMinContactLegs} legs must
remain in ground contact, so at most
\numberstringnum{\corlMaxConcurrentSwingLegs} legs may swing at once.
This constraint gives the robot stability margin. Its effect at the
easiest and hardest terrain settings has not yet been ablated.
\changed{The actor is also trained on
\numberstringnum{\corlTrainCandidatesPerLeg} randomly sampled
candidates per leg. At deployment, it is evaluated on an exhaustive
enumeration of the same discrete terrain grid. The deployed candidate
set is a denser version of the training set, not a different action
type. Still, we have not measured performance as a function of
candidate-set size. That study would separate the benefit of dense
selection from the benefit of the learned value function itself.}

Both simulation and hardware rely on an external terrain scan to build
the validity mask. Simulation uses a raycast height field. The
hardware system uses the Legged Perceptive Stack's height map. This
difference may create sim-to-real error, but we have not quantified
it. Future work will add direct exteroceptive input, such as point
clouds or height maps processed by convolutional layers. This should
make the planner less dependent on hand-tuned terrain thresholding. We
also plan to expand the action beyond foothold and timing. For
example, the planner could choose body posture or center-of-mass
targets. This would let the same high-rate re-decision idea address
steeper terrain and stairs without hand-coded elevation heuristics.
